% ============================================================================
% chap5_safety_eval.tex
%   第5章 §5.2.3 · 安全评估与模型验证机制
%   标准化任务配置 → 策略执行 (内建/上传) → 统一指标体系 → 对比报表
%
% 独立编译:
%   xelatex figs/chap5_safety_eval.tex
%   → figs/chap5_safety_eval.pdf
%
% 嵌入论文:
%   \begin{figure}[htbp]
%     \centering
%     \includegraphics[width=0.95\linewidth]{figs/chap5_safety_eval.pdf}
%     \caption{安全评估与模型验证机制: 标准化任务配置统一输入, 内建算法与用户
%              上传模型在同一环境下执行, 统一指标体系 (任务性能 + 安全性)
%              汇总至对比报表, 支持跨算法横向比较.}
%     \label{fig:chap5_safety_eval}
%   \end{figure}
% ============================================================================

\documentclass[tikz,border=6pt]{standalone}
\usepackage[UTF8]{ctex}
\usepackage{amsmath}
\usepackage{tikz}
\usetikzlibrary{positioning, arrows.meta, shapes.geometric, calc, fit, backgrounds}

% 配色与前几张图保持一致：文字统一黑色，橙色边框统一为深橙色
\definecolor{cOurs}{HTML}{0B4F8A}      % 深蓝：策略来源
\definecolor{cOursAlt}{HTML}{C05600}   % 深橙：配置与报表
\definecolor{cOursEval}{HTML}{000000}  % 黑色：评估流水线与指标体系
\definecolor{cUpstream}{HTML}{4D4D4D}  % 深灰：通用基础模块
\definecolor{cBgLayer}{HTML}{F0F0F0}
\definecolor{cBgHi}{HTML}{FFF8E7}

\begin{document}
\begin{tikzpicture}[
  font=\large\sffamily,
  every node/.style={align=center, text=black},
  % 标准化配置源：橙色框，字体略小
  configbox/.style={
    draw=cOursAlt, line width=1.1pt, rounded corners=3pt,
    fill=cBgHi, text=black, inner sep=4pt,
    minimum width=4.2cm, minimum height=2.2cm,
    font=\small\bfseries\sffamily,
  },
  % 策略来源：蓝色框，字体略小
  srcbuiltin/.style={
    draw=cOurs, line width=1pt, rounded corners=3pt,
    fill=white, text=black, inner sep=4pt,
    minimum width=3.6cm, minimum height=1.05cm,
    font=\small\sffamily,
  },
  srcupload/.style={
    draw=cOurs, line width=1pt, rounded corners=3pt,
    fill=cOurs!5, text=black, inner sep=4pt,
    minimum width=3.6cm, minimum height=1.05cm,
    font=\small\sffamily,
  },
  % 验证流水线阶段：黑色框，字体放大
  stage/.style={
    draw=cOursEval, line width=1pt, rounded corners=3pt,
    fill=white, text=black, inner sep=3pt,
    minimum width=2.3cm, minimum height=1.25cm,
    font=\small\sffamily,
  },
  % 指标体系容器：黑色框，字体放大
  metricbox/.style={
    draw=cOursEval, line width=1pt, rounded corners=3pt,
    fill=white, inner sep=4pt,
    minimum width=3.3cm, minimum height=2.6cm,
    font=\small\sffamily, text=black,
  },
  metricinner/.style={
    draw=cOursEval!60, line width=0.4pt, rounded corners=1.5pt,
    fill=cOursEval!6, text=black, inner sep=1.5pt,
    minimum width=2.85cm, minimum height=0.36cm,
    font=\footnotesize\sffamily,
  },
  mtitle/.style={
    font=\small\bfseries\sffamily, text=black,
    inner sep=1pt,
  },
  % 输出报表：橙色框，字体略小
  report/.style={
    draw=cOursAlt, line width=1.1pt, rounded corners=3pt,
    fill=cBgHi, text=black, inner sep=4pt,
    minimum width=3.6cm, minimum height=1.8cm,
    font=\small\bfseries\sffamily,
  },
  % 箭头
  flow/.style={
    -{Stealth[length=2.4mm, width=2mm]},
    line width=0.9pt, draw=black!70,
  },
  flowCfg/.style={
    -{Stealth[length=2.2mm, width=1.8mm]},
    line width=0.9pt, draw=cOursAlt,
  },
  flowPolicy/.style={
    -{Stealth[length=2.2mm, width=1.8mm]},
    line width=0.9pt, draw=cOurs,
  },
  flowMetric/.style={
    -{Stealth[length=2.2mm, width=1.8mm]},
    line width=0.9pt, draw=cOursEval,
  },
  lbl/.style={font=\footnotesize, inner sep=1pt, fill=white, text=black},
]

% ============================================================================
% 左侧: 标准化任务配置 (输入源, 橙)
% ============================================================================
\node[configbox] (cfg) at (-7.5cm, 0)
  {标准化任务配置\\[3pt]
   \footnotesize 机器人模型 · 场景结构\\
   \footnotesize 障碍物分布 · 传感器参数\\
   \footnotesize 随机种子 \& 评估轮次};

% ============================================================================
% 左下: 策略来源 (内建 + 用户上传, 蓝)
% ============================================================================
\node[srcbuiltin] (builtin) at (-7.5cm, -2.8cm)
  {内建算法库\\[-1pt]
   \footnotesize MAPPO + 安全机制};

\node[srcupload, below=0.35cm of builtin] (upload)
  {用户上传模型\\[-1pt]
   \footnotesize .pt / .onnx 检查点上传接口};

% ============================================================================
% 中部: 统一验证流水线 (黑色, 4 阶段)
% ============================================================================
\node[stage] (s1) at (-3.3cm, -0.0cm)
  {环境初始化\\[-1pt]
   \footnotesize 加载 Task 配置\\[-1pt]
   \footnotesize 同一 seed 流};

\node[stage] (s2) at (-0.55cm, -0.0cm)
  {策略推理\\[-1pt]
   \footnotesize 固定 episode \&\\[-1pt]
   \footnotesize 控制频率};

\node[stage] (s3) at (2.2cm, -0.0cm)
  {运行记录\\[-1pt]
   \footnotesize traj / cost /\\[-1pt]
   \footnotesize constraint 日志};

\node[stage] (s4) at (4.95cm, -0.0cm)
  {指标计算\\[-1pt]
   \footnotesize 聚合 \&\\[-1pt]
   \footnotesize 均值/方差};

\draw[flowMetric] (s1.east) -- (s2.west);
\draw[flowMetric] (s2.east) -- (s3.west);
\draw[flowMetric] (s3.east) -- (s4.west);

% 流水线横幅
\node[font=\small\bfseries\sffamily, text=black,
      fill=cOursEval!6, draw=cOursEval!50, rounded corners=2pt, inner sep=3pt]
  at (0.85cm, 1.35cm)
  {统一验证流水线 (同一环境 · 同一随机种子 · 同一评估轮次)};

% 配置 → 流水线 (橙线)
\draw[flowCfg] (cfg.east) to[out=0,in=180]
  node[lbl, pos=0.55]{task config}
  (s1.west);

% 策略 → 推理阶段 (蓝线)
\draw[flowPolicy] (builtin.east) to[out=0,in=-120]
  node[lbl, pos=0.6]{policy $\pi$}
  (s2.south);

\draw[flowPolicy] (upload.east) to[out=0,in=-90]
  ($(s2.south)+(0.25cm,0)$);

% ============================================================================
% 右侧: 统一指标体系 (两列: 性能 / 安全)
% ============================================================================
\node[metricbox] (mperf) at (8.2cm, 1.55cm) {};
\node[mtitle, below=0.08cm of mperf.north, anchor=north] (mperft)
  {任务性能指标};

\node[metricinner, below=0.08cm of mperft.south, anchor=north] (mp1)
  {任务完成率};
\node[metricinner, below=0.08cm of mp1.south, anchor=north] (mp2)
  {累计回报 EpRet};
\node[metricinner, below=0.08cm of mp2.south, anchor=north] (mp3)
  {平均编队误差};
\node[metricinner, below=0.08cm of mp3.south, anchor=north] (mp4)
  {终点距目标 $d_\mathrm{final}$};

\node[metricbox] (msafe) at (8.2cm, -1.55cm) {};
\node[mtitle, below=0.08cm of msafe.north, anchor=north] (msafet)
  {安全性指标};

\node[metricinner, below=0.08cm of msafet.south, anchor=north] (ms1)
  {约束违反次数};
\node[metricinner, below=0.08cm of ms1.south, anchor=north] (ms2)
  {最小两两距离};
\node[metricinner, below=0.08cm of ms2.south, anchor=north] (ms3)
  {最小墙距};
\node[metricinner, below=0.08cm of ms3.south, anchor=north] (ms4)
  {累计安全代价 $\sum c_t$};

% 流水线输出 → 两类指标
\draw[flowMetric] ($(s4.east)+(0,0.14cm)$) to[out=10,in=180]
  node[lbl, pos=0.5]{metrics}
  ($(mperf.west)+(0,-0.12cm)$);

\draw[flowMetric] ($(s4.east)+(0,-0.14cm)$) to[out=-10,in=180]
  ($(msafe.west)+(0,0.12cm)$);

% ============================================================================
% 右下: 对比报表 (橙)
% ============================================================================
\node[report] (rep) at (8.2cm, -4.4cm)
  {跨算法对比报表\\[2pt]
   \footnotesize 表格 · 轨迹 · 学习曲线\\
   \footnotesize 一键导出 (CSV / PDF)};

   \coordinate (repR1) at ($(rep.east)+(0,0.35cm)$);
   \coordinate (repR2) at ($(rep.east)+(0,-0.35cm)$);
   \coordinate (rightBus) at ($(rep.east)+(0.65cm,0)$);
   
   \draw[flowMetric] (mperf.east) -- ($(mperf.east)+(0.65cm,0)$)
                     |- (repR1);
   
   \draw[flowMetric] (msafe.east) -- ($(msafe.east)+(0.65cm,0)$)
                     |- (repR2);

% ============================================================================
% 左下: 图例
% ============================================================================
\node[draw=black!25, rounded corners=2pt, inner sep=5pt, font=\small,
      text=black, anchor=north west]
  at (-9.0cm, -5.2cm)
  {\begin{tabular}{@{}l@{\hspace{4pt}}l@{\hspace{14pt}}l@{\hspace{4pt}}l@{\hspace{14pt}}l@{\hspace{4pt}}l@{}}
     \tikz[baseline=-0.6ex]\draw[flowCfg] (0,0) -- (0.6,0); & 配置下发 &
     \tikz[baseline=-0.6ex]\draw[flowPolicy] (0,0) -- (0.6,0); & 策略注入 &
     \tikz[baseline=-0.6ex]\draw[flowMetric] (0,0) -- (0.6,0); & 评估数据流 \\
   \end{tabular}};

\end{tikzpicture}
\end{document}